% ============================================================
\section{Experimental Setup}
\label{sec:experiments}
% ============================================================

We evaluate routing algorithms through cycle-accurate simulations using BookSim2~\cite{booksim}, a widely-used NoC simulator. This section describes the experimental methodology, baselines, traffic patterns, and evaluation metrics.

\subsection{Simulation Environment}

Simulation parameters are summarized in Table~\ref{tab:simparams}. All results are obtained with 100 independent sample periods of 1,000 cycles each, preceded by 3 warmup periods.

\begin{table}[h]
\centering
\caption{BookSim2 simulation parameters}
\label{tab:simparams}
\begin{tabular}{ll}
\toprule
\textbf{Parameter} & \textbf{Value} \\
\midrule
Simulator & BookSim2 2.0 \\
Topologies & Mesh $4\times4$ \\
Flit size & 32 bytes \\
Buffer depth & 4 flits/VC \\
Virtual channels & 4 \\
Routing computation & 1 cycle \\
Warmup periods & 3 \\
Sample periods & 100 (1,000 cycles each) \\
Simulation length & $\sim$103,000 cycles \\
Packet size & 1 flit \\
\midrule
Seeds per config & 5 \\
Total runs & 450 (across 5 algorithms, 3 traffics, 6 rates, 5 seeds) \\
\bottomrule
\end{tabular}
\end{table}

\subsection{Routing Algorithms}

We compare five routing approaches:

\begin{itemize}
    \item \textbf{XY (DOR):} Dimension-order routing, the most widely used deterministic algorithm for mesh NoCs. Packets are routed first along the X dimension, then Y.
    \item \textbf{Adaptive XY/YX:} An adaptive variant that selects between XY and YX paths based on local congestion at the source router~\cite{balfour2006}.
    \item \textbf{Minimal Adaptive:} Routes packets along any minimal path, with port selection based on buffer occupancy.
    \item \textbf{Valiant:} Randomized routing that sends packets to a random intermediate node before forwarding to the destination.
    \item \textbf{GNNocRoute-PPO:} Precomputed routing policy based on a GATv2-GNN encoder (see Section~\ref{sec:method}). The model produces per-source-destination routing decisions (XY or YX) that are stored in a lookup table for online use.
\end{itemize}

\subsection{Traffic Patterns}

We evaluate under three synthetic traffic patterns:

\begin{itemize}
    \item \textbf{Uniform Random:} Each node sends packets to a uniformly random destination. Represents balanced workloads.
    \item \textbf{Transpose:} Node $(i,j)$ sends to node $(j,i)$. Creates diagonal traffic that stresses non-minimal path selection.
    \item \textbf{Hotspot (10\%):} $10\%$ of traffic targets a single central node $(3,3)$, with the remaining $90\%$ uniformly distributed. Tests congestion handling under skewed loads.
\end{itemize}

\subsection{Evaluation Metrics}

We measure two primary metrics:

\begin{itemize}
    \item \textbf{Average packet latency:} Mean number of cycles from packet injection to tail flit reception, averaged over 100 sample periods and 5 random seeds.
    \item \textbf{Saturation behavior:} Injection rate at which latency exceeds the simulation stability threshold (500 cycles), indicating network saturation.
\end{itemize}

\subsection{DRL Training Setup}

The GNN-PPO routing policy is trained offline using the NoC routing environment (Section~\ref{sec:method}). Training is conducted on Mesh $4\times4$ with 500 episodes of GATv2 encoder pre-training followed by 500 episodes of DQN training. The trained policy is exported as a precomputed routing table.

All training runs on CPU (Intel Xeon, 2.3 GHz) using PyTorch 2.8 and PyTorch Geometric 2.6.
